%A compiler en XELATEX sinon pas d'arbre de proba
\documentclass[french]{book}
\usepackage{teach}
\usepackage{verbatim}
\usepackage[french]{babel}
%%% Couleurs
\redefinecolor{thm}{title}{bgcolor}{alizarin}
\redefinecolor{prop}{title}{bgcolor}{meatbrown}
\redefinecolor{defn}{title}{bgcolor}{olivine}
\definecolor{alizarin}{rgb}{0.82, 0.1, 0.26}
\definecolor{meatbrown}{rgb}{0.9, 0.72, 0.23}
\definecolor{olivine}{rgb}{0.6, 0.73, 0.45}
\usepackage{color}
\usepackage[table]{xcolor} 
%%%%%%Les symboles
\usepackage{amsmath}
\usepackage{amssymb}
\usepackage{amsfonts}
\usepackage{amsthm}
\usepackage{mwe}
\usepackage{yhmath} 
\usepackage{eurosym}
%%% Ensemble de nombre
\newcommand{\R}{\mathbb{R}}
%\newcommand{\C}{\mathds {C}}
\newcommand{\N}{\mathbb{N}}
\newcommand{\Z}{\mathbb{Z}}
\newcommand{\Q}{\mathbb{Q}}
\newcommand{\D}{\mathbb{D}}
%%% Tableau
\usepackage{array}
\usepackage{multicol}
\setlength{\multicolsep}{3pt}
\usepackage{tkz-tab}
\usepackage{cellspace}
\usepackage{diagbox}
%%% Figures
\usepackage{tikz}
\usepackage{graphicx}
\graphicspath{{Figures/}}
%\usepackage{pstricks,pst-plot,pst-text,pst-tree,pst-eucl}
%%%Affichage
\newcommand{\ds}{\displaystyle}
\renewcommand{\arraystretch}{1.5} 
\newcolumntype{C}[1]{>{\centering\arraybackslash}p{#1cm}}
%%% Lecon creuse/pleine
\usepackage{dashundergaps}
%\TeacherModeOff
\TeacherModeOn
%%% Macro
\newcommand{\V}{\overrightarrow}
\newcommand{\Rep}{(O;\V{i};\V{j})}
\newcommand{\Syst}[2]{\left\{\begin{array}{lcccc} #1 \\ #2 \end{array}\right.}
\begin{document}
\setcounter{chapter}{8}
\chapter{Fonction inverse}
\vspace*{-1cm}
%La fonction inverse est particulière, on la retrouve géométriquement dans l'application du théorème de Thalès.\\
%La fonction inverse trouve son utilité dans la construction de nombre fractionnaire, à tout nombre réel strictement positif $x$ on peut placer géométriquement $\dfrac{1}{x}$ à l'aide du Théorème de Thalès ce qui correspond à un changement d'échelle.
%
%Autrement dit: "\textit{Si un nombre $x$ \textbf{non nul} devient l'unité et que l'écart est respecté de manière proportionnelle alors que devient $1$ ?}" Il deviendra en réalité exactement $\dfrac{1}{x}$.

\section{Étude de la fonction inverse}


\begin{defn}
On appelle fonction inverse la fonction définie pour tout nombre réel appartenant à $]-\infty ; 0[ \cup ] 0 ;+\infty[$ \emph{(aussi notée $\mathbb{R}^*$)} par $f(x)=\dfrac{1}{x}$.
\end{defn}


\begin{exemple}[Tableau de valeurs]
\begin{center}
\renewcommand{\arraystretch}{2.3}
\begin{tabular}	{|Sl||Sl|Sl|Sl|Sl|Sl|Sl|Sl|Sl|Sl|Sl|}
\hline  $x$ & -4 & -3 & -2 & -1 & $-\dfrac{1}{2}$ & 0 & $\dfrac{1}{2}$ & 1 & 2 & 4 \\
\hline $\dfrac{1}{x}$ & $-\dfrac{1}{4}$ & - $\dfrac{1}{3}$ & - $\dfrac{1}{2}$ & $-1$ & -$2$ & $\times$ & $2$ & $1$ & $\dfrac{1}{2}$ & $\dfrac{1}{4}$ \\
\hline
\end{tabular}
\end{center}
\end{exemple}

\begin{rem}
La fonction inverse est une \gap*{involution} sur $\mathbb{R}^*$, autrement dit, l'inverse de l'inverse d'un nombre non nul est lui-même.
\end{rem} 

\section{Variations, représentation graphique et signe}

\subsection{Variations}

%%%%% V1 - Avec Dérivation+preuve%%%%
\begin{prop}
La fonction inverse est dérivable sur $]-\infty ; 0[$ et sur $] 0 ;+\infty[$ et sa dérivée est $-\dfrac{1}{x^2}$
\end{prop}
 
\begin{demo}
On a pour tout réel $a$ non nul et tout réel $h \neq 0$ tel que $a+h \neq 0$, 
$$\dfrac{f(a+h)-f(a)}{h}=\dfrac{\dfrac{1}{a+h}-\dfrac{1}{a}}{h}=\dfrac{\dfrac{a-(a+h)}{a(a+h)}}{h}=\dfrac{\dfrac{-h}{a(a+h)}}{h}=-\dfrac{1}{a(a+h)}$$

Par suite en faisant tendre $h$ vers $0$, on obtient $$f'(a):=\lim_{h\mapsto0} \dfrac{f(a+h)-f(a)}{h}=-\dfrac{1}{a^2}$$
\end{demo}

\begin{prop}
\begin{enumerate}
\item La fonction inverse est \gap*{strictement décroissante} sur $\mathbb{R}_+^*$ car sa dérivée est \gap*{strictement négative} sur $\mathbb{R}_+^*$.

\item La fonction inverse est \gap*{strictement décroissante} sur $\mathbb{R}_-^*$ car sa dérivée est \gap*{strictement négative} sur $\mathbb{R}_-^*$.
\end{enumerate}
\end{prop}

\begin{attention}
Il est \textbf{faux} de dire que la fonction inverse est strictement décroissante sur $\mathbb{R}^*$. Le contre-exemple suivant suffit à comprendre pourquoi. 
$$-2<2 \;\; \textrm{et pourtant,} \;\; -\dfrac{1}{2} < \dfrac{1}{2}$$ 
\end{attention}

%%%%%%%%%%%%%%%%%

%%%%% V2 - Sans Dérivation+preuve%%%%
%\begin{prop}
%La fonction inverse est \gap*{strictement décroissante} sur $\mathbb{R}^*_-$ et  \gap*{strictement décroissante} sur $\mathbb{R}^*_+$.
%\end{prop}
%
%\begin{attention}
%Il est \textbf{faux} de dire que la fonction inverse est strictement décroissante sur $\mathbb{R}^*$. Le contre-exemple suivant suffit à comprendre pourquoi. 
%$$-2<2 \;\; \textrm{et pourtant,} \;\; -\dfrac{1}{2} < \dfrac{1}{2}$$ 
%\end{attention}
%
%\begin{demo}
%
%\underline{Preuve de la stricte décroissance sur $\mathbb{R}^*_-$ de la fonction inverse:}
%
%Soit $a<b < 0$ alors $\frac{1}{a} - \frac{1}{b} = \frac{b-a}{ab} $ , ici $ab > 0$ car $a<b<0$ et $0<b-a$ car $a<b$. Donc $\frac{1}{a} - \frac{1}{b} >0$. 
%$$\textrm{D'où} \;\; \dfrac{1}{a}>\dfrac{1}{b}$$
%
%\underline{Preuve de la stricte décroissance sur $\mathbb{R}^*_+$ de la fonction inverse:}
%
%Soit $0<a<b$ alors $\frac{1}{a} - \frac{1}{b} = \frac{b-a}{ab} $ , ici $ab > 0$ car $0<a<b$ et $0<b-a$ car $a<b$. Donc $\frac{1}{a} - \frac{1}{b} >0$. 
%$$\textrm{D'où} \;\; \dfrac{1}{a}>\dfrac{1}{b}$$
%\end{demo}
%%%
\subsection{Représentation graphique}

\begin{prop}
La représentation graphique de la fonction inverse est appelée \gap*{hyperbole}, ceci est dû à sa forme aux bords de son domaine de définition.
\end{prop}


\vspace*{-6cm}
\begin{center}
\hspace*{8.5cm}\includegraphics[scale=0.85]{inverse.1}
\end{center}

\vspace*{5.7cm}
\begin{rem}
La courbe de la fonction inverse possède \gap*{un centre de symétrie}, à savoir l'origine du repère. On dit aussi que la fonction est \gap*{impaire}, autrement dit  $f(-x)=-f(x)$.
\end{rem}

\begin{defn}
Une \gap*{asymptote} est une droite dont une courbe s'approche de plus en plus, sans jamais l'atteindre.
\vspace*{5mm}
\end{defn}


\begin{rem}[Comportement aux bornes de son ensemble de définition]

La représentation graphique admet deux \gap*{asymptotes}: l'axe des ordonnées et l'axe des abscisses. 

\begin{itemize}
\item Lorsque $x$ devient de plus en plus proche de $0$,  $\dfrac{1}{x}$ devient en valeur absolue très grand et tend vers $+\infty$. Autrement dit, la courbe se rapproche de l'axe des ordonnées sans jamais la rencontrer, on dit que c'est une \gap*{asymptote verticale}  à la courbe $\mathcal{C}_f$.
\item Lorsque $x$ devient de plus en plus grand en valeur absolue, $\dfrac{1}{x}$ tend alors vers $0$. Autrement dit, la courbe se rapproche de l'axe des abscisses sans jamais la rencontrer, on dit que c'est une \gap*{ asymptote horizontale} à la courbe $\mathcal{C}_f$.
\end{itemize}
\end{rem}

\subsection{Signe}

\begin{prop}
La fonction inverse \gap*{conserve le signe}, elle est donc \gap*{strictement négative} sur $\mathbb{R}_-^*$ et \gap*{strictement positive} sur $\mathbb{R}_+^*$.
\end{prop}

\section{Application à l'étude de fonctions}

%\begin{exemple}[Étude d'une fonction somme d'une fonction polynomiale et de la fonction inverse]
%
%On considère la fonction $f$ définie sur $] 0 ;+\infty[$ par $ f(x)=2 x+1+\dfrac{8}{x}$.
%
%On a $f^{\prime}(x)=2-\dfrac{8}{x^2}$. Comme pour tout réel $x$ non nul, $x^{2}>0$, $ f^{\prime}(x)$ est du signe de $2x^2-8$ qui est une fonction polynôme du second degré. Par ailleurs, on constate que $f'(x)=2(x^2-4)=2(x+2)(x-2)$. De plus $a>0$ d'où le tableau récapitulatif suivant : \\
%\begin{center}
%\begin{tikzpicture}
%     \tkzTabInit[lgt=3,espcl=3]
%       {$x$ / 1 , Signe de $f'(x)$ / 1, Variation de  $f$ / 2 }
%       {$0$, 2 , $+\infty$}
%       \tkzTabLine
%     		{ d , - , 0 , + , }
%       \tkzTabVar
%     		{ D+/ , -/$f(2)=9$ , +/$0$}
%\end{tikzpicture}
%\end{center}
%\end{exemple}

\begin{memento}
Une application de la fonction inverse pour une entreprise est l'étude des coûts moyens de production. Si l'on note $C_T(q)$ le \textbf{coût total de production} en fonction du nombre d'objet vendu noté $q$. Alors par définition on a le \textbf{coût moyen de production}
$$C_m(q):= \dfrac{C_T(q)}{q}$$ 
Le but pour une entreprise est de \gap*{minimiser} cette quantité on fonction de $q$. Cependant attention, $q$ n'est pas un nombre réel quelconque mais un \gap*{nombre strictement positif} \textit{(entier ou non en fonction de l'objet de la production)} qui ne doit pas dépasser la capacité de production maximale $q_{max}$ de l'entreprise. 
\end{memento}

\newpage
\section{Exercices}
%%%%Edupuy

\begin{exocor}
Résoudre les équations et inéquations suivantes :
\begin{multicols}{3}
\begin{enumerate}
\item $2+\dfrac{1}{x}=3$
\item $\dfrac{2}{x}=6$
\item $1+\dfrac{1}{x} \geq 2 $
\item $-6+\dfrac{2}{x}=\dfrac{1}{x}+2$
\item $1+\dfrac{1}{x} > 2 $
\item $1+\dfrac{1}{x} < 1+ 10^{-3}  $
\end{enumerate}
\end{multicols}
\end{exocor}

\begin{exocor}
Sachant que $1<x\leq 3$ encadrer les quantités suivantes :
\begin{multicols}{4}
\begin{enumerate}
\item $\dfrac{1}{x}$
\item $\dfrac{2}{x}$
\item $1+\dfrac{1}{x} $
\item $-6+\dfrac{2}{x}$
\end{enumerate}
\end{multicols}
\end{exocor}

\begin{exocor}
Soit $x \neq -2$ un nombre réel. 
\begin{enumerate}
\item Pour quelles valeurs de $x$ le réel $A=1- \dfrac{2}{x+2}$ admet-il un inverse ?
\item Donner l'expression en fonction de $x$ de l'inverse de $A$.
\end{enumerate}
\end{exocor}

\begin{exocor}
On note  $f$ la fonction inverse :
\begin{enumerate}
	\item Quel est l'ensemble de définition de la fonction $f$ ?
	\item Calculer l'image par $f$ de chacun des réels suivants : $(- 0,25)$ ; $10^{-3}$ ; $\left(-\dfrac{3}{4}\right)$  et $1$.
	\item Déterminer le réel dont l'image par $f$ est $(-10)$, puis celui dont l'image par $f$ est $0,01$.
\end{enumerate}
\end{exocor}


\begin{exocor}
Calculer les dérivées des fonctions suivantes sur $\mathbb{R}^*$.

\begin{multicols}{3}
\begin{enumerate}
\item $f(x)= x+\dfrac{1}{x}$
\item $g(x)=\dfrac{2}{x}$
\item $h(x)=-x^2-\dfrac{2}{x}$
\item $i(t)=3t-2+\dfrac{1}{t}$
\item $j(x)=2x+\dfrac{3}{x}$
\item $k(x)=7x^2-2x+6+\dfrac{4}{x}$
\item $l(x)=x^2 + \dfrac{1}{x}$
\item $m(x)=3x^2-5x+2+\dfrac{1}{x}$
\item $n(x)=-2x^3+4x^2+1+\dfrac{1}{x}$
\end{enumerate}
\end{multicols}
\end{exocor}

\begin{exocor}
Étudier le sens de variations des fonctions suivantes, définies sur $\mathbb{R}^*$.

\begin{multicols}{2}
\begin{enumerate}
\item $f(x)=3+\dfrac{1}{x}$
\item $g(x)=2-\dfrac{1}{x}$
\end{enumerate}
\end{multicols}
\end{exocor}


\begin{exocor}
Soit $f$ la fonction définie sur $\mathbb{R}^*$ par : $f(x)=x+\dfrac{1}{x}$

\begin{enumerate}
\item Calculer $f'(x)$.
\item Montrer que pour tout réel $x$ non nul : $f'(x)=\dfrac{(x-1)(x+1)}{x^2}$.
\item Après avoir étudié le signe de $f'(x)$, dresser le tableau de variations de $f$.
\end{enumerate}
\end{exocor}

\begin{exocor}
Soit $f$ la fonction définie sur $\mathbb{R}^*$ par : $f(x)=4x+1+\dfrac{9}{x}$

\begin{enumerate}
\item Calculer $f'(x)$.
\item Montrer que pour tout réel $x$ non nul : $f'(x)=\dfrac{(2x-3)(2x+3)}{x^2}$.
\item Après avoir étudié le signe de $f'(x)$, dresser le tableau de variations de $f$.
\end{enumerate}
\end{exocor}

\begin{exocor}
Soit $g$ la fonction définie sur $\mathbb{R}^*$ par : $g(x)=x+10-\dfrac{1}{x}$

\begin{enumerate}
\item Calculer $g'(x)$.
\item Montrer que pour tout réel $x$ non nul : $g'(x)=\dfrac{x^2+1}{x^2}$.
\item Après avoir étudié le signe de $g'(x)$, dresser le tableau de variations de $g$.
\end{enumerate}
\end{exocor}

\begin{exocor}
Soit $h$ la fonction définie sur $\mathbb{R}^*$ par : $h(x)=\dfrac{3}{2}x^2-7x-\dfrac{4}{x}$

\begin{enumerate}
\item Calculer $h'(x)$.
\item Montrer que pour tout réel $x$ non nul : $h'(x)=\dfrac{(x-2)(x-1)(3x+2)}{x^2}$.
\item Après avoir étudié le signe de $h'(x)$, dresser le tableau de variations de $h$.
\end{enumerate}
\end{exocor}

\vfill

\begin{exocor}
Soit $f$ la fonction définie sur l'intervalle $[40 ; 160]$ par : 
$f(x) = x-100+ \dfrac{6 400}{x}$.\\

La courbe ci-dessous est la représentation graphique de f dans un repère.\\
  %%% representation graph
\hspace*{2cm}\includegraphics[scale=1]{inverse.2}\\

\begin{enumerate}
\item Calculer la dérivée $f'$ de $f$ et montrer qu'elle peut s'écrire : $f'(x) = \dfrac{x^2 - 6 400}{x^2}$.
\item Étudier le signe de $f'(x)$ sur l'intervalle $[40 ; 160]$.
\item En déduire le tableau des variations de $f$ sur l'intervalle $[40 ; 160]$.
\item Le coût, exprimé en euros, de $x$ repas préparés par service dans un restaurant, peut s'écrire, pour $x\in[40; 160]$:
$$C(x) = x^2-100x +6 400$$
\begin{enumerate}
\item Recopier et compléter le tableau:
\begin{center}
\begin{tabular}{|c|c|c|c|}
\hline 
Nombre de repas & 40 & 50 & 100 \\ 
\hline 
Coût de $x$ repas & • & • & • \\ 
\hline 
 Coût moyen d'un repas & • & • & • \\ 
\hline 
\end{tabular} 
\end{center}
\item Écrire le coût moyen d'un repas en fonction du nombre $x$ de repas préparés. On notera ce coût moyen unitaire $C_m(x)$.
\item  Déduire de la question 3,le nombre de repas que ce restaurant doit servir pour que le coût moyen d'un repas soit minimal.
\item Trouver, à l'aide du graphique, à quel intervalle doit appartenir $x$ pour que le coût moyen unitaire soit inférieur ou égal à $90$ euros.
\end{enumerate}
\end{enumerate}
\end{exocor}

\newpage
\section{Corrigés}

\begin{corrige}
\begin{multicols}{3}
\begin{enumerate}
\item $x=1$
\item $x=\dfrac{1}{3}$
\item $x \leq 1$
\item $x=\dfrac{1}{8}$
\item $x<1$
\item $x>10^3$
\end{enumerate}
\end{multicols}
\end{corrige}

\begin{corrige}
\begin{multicols}{4}
\begin{enumerate}
\item $1 > \dfrac{1}{x} \geq \dfrac{1}{3}$
\item $2 > \dfrac{2}{x} \geq \dfrac{2}{3}$
\item $2 >1+\dfrac{1}{x} \geq \dfrac{4}{3} $
\item $-4 >-6+\dfrac{2}{x} \geq \dfrac{-16}{3}$
\end{enumerate}
\end{multicols}
\end{corrige}

\begin{corrige}
\begin{enumerate}
\item $A$ n'admet pas d'inverse si et seulement si $A=0$. 
$$A=0 \Longleftrightarrow 1-\dfrac{2}{x+2} =0 \Longleftrightarrow 1=\dfrac{2}{x+2} \Longleftrightarrow 2=x+2 \Longleftrightarrow x=0$$
Autrement dit $A$ n'admet pas d'inverse si $x=0$. 
\item Supposons que $x \neq -2$ et $x\neq 0$ alors 
$$A=1-\dfrac{2}{x+2}=\dfrac{x}{x+2}$$
Donc 
$$\dfrac{1}{A}=\dfrac{1}{\frac{x}{x+2}}=\dfrac{x+2}{x}$$

\end{enumerate}
\end{corrige}

\begin{corrige}
\begin{enumerate}
\item L'ensemble de définition de $f$ est  $\mathbb{R}^*$.
\item $f(-0,25)=-4$ ; $f(10^{-3})=10^3$ ; $(-\dfrac{3}{4})$ ; $f(1)=1$.
\item $f(-0,1)=-10$ et $f(100)=0,01$
\end{enumerate}
\end{corrige}

\begin{corrige}
\begin{multicols}{3}
\begin{enumerate}
\item $f'(x)= 1-\dfrac{1}{x^2}$
\item $g'(x)=-\dfrac{2}{x^2}$
\item $h'(x)=-2x+\dfrac{2}{x^2}$
\item $i'(t)=3-\dfrac{1}{t^2}$
\item $j'(x)=2-\dfrac{3}{x^2}$
\item $k'(x)=14x-2-\dfrac{4}{x^2}$
\item $l'(x)=2x-\dfrac{1}{x^2}$
\item $m'(x)=6x-5-\dfrac{1}{x^2}$
\item $n'(x)=-6x^2+8x-\dfrac{1}{x^2}$
\end{enumerate}
\end{multicols}
\end{corrige}

\begin{corrige}
\begin{enumerate}
\item On sait que $f'(x)=-\dfrac{1}{x^2}$ qui est strictement négative sur $\mathbb{R}^*$ donc $f$ est strictement décroissante sur chacun des intervalles de $\mathbb{R}^*$.

\begin{center}
\begin{tikzpicture}
     \tkzTabInit[lgt=3,espcl=3]
       {$x$ / 1 , Signe de $f'(x)$ / 1, Variation de  $f$ / 2 }
       {$-\infty$,  0 , $+\infty$}
       \tkzTabLine
     		{  , - , d , - , }
       \tkzTabVar
     		{ +/ , -D+/$ $/ $ $ , -/  }
\end{tikzpicture}
\end{center}

\item  On sait que $f'(x)=\dfrac{1}{x^2}$ qui est strictement positive sur $\mathbb{R}^*$ donc $f$ est strictement croissante sur chacun des intervalles de $\mathbb{R}^*$.

\begin{center}
\begin{tikzpicture}
     \tkzTabInit[lgt=3,espcl=3]
       {$x$ / 1 , Signe de $f'(x)$ / 1, Variation de  $f$ / 2 }
       {$-\infty$,  0 , $+\infty$}
       \tkzTabLine
     		{  , + , d , + , }
       \tkzTabVar
     		{ -/ , +D-/$ $/ $ $ , +/  }
\end{tikzpicture}
\end{center}
\end{enumerate}
\end{corrige}


\begin{corrige}
\begin{enumerate}

\item $f'(x)=1-\dfrac{1}{x^2}$
\item D'une part on sait que $f'(x)=1-\dfrac{1}{x^2}=\dfrac{x^2-1}{x^2}$, d'autre part $(x-1)(x+1)=x^2-1^2= x^2-1$ d'où $f'(x)=\dfrac{(x-1)(x+1)}{x^2}$.
\item Tableau récapitulatif d'étude de la fonction $f$,
\begin{center}
\begin{tikzpicture}
     \tkzTabInit[lgt=3,espcl=3]
       {$x$ / 1 , Signe de $f'(x)$ / 1, Variation de  $f$ / 2 }
       {$-\infty$, -1 ,  0 , 1,  $+\infty$}
       \tkzTabLine
     		{ ,+,z,-,d,-,z,+, }
       \tkzTabVar
     		{ -/ , +/ $-2$, -D+ / / ,-/ $2$, +/   }
\end{tikzpicture}
\end{center}
\end{enumerate}
\end{corrige}

\begin{corrige}
\begin{enumerate}

\item $f'(x)=4-\dfrac{9}{x^2}$
\item D'une part on sait que $f'(x)=4-\dfrac{9}{x^2}=\dfrac{4x^2-9}{x^2}$, d'autre part $(2x-3)(2x+3)=(2x)^2-3^2= 4x^2-9$ d'où $f'(x)=\dfrac{(2x-3)(2x+3)}{x^2}$.
\item Tableau récapitulatif d'étude de la fonction $f$,
\begin{center}
\begin{tikzpicture}
     \tkzTabInit[lgt=3,espcl=3]
       {$x$ / 1 , Signe de $f'(x)$ / 1, Variation de  $f$ / 2 }
       {$-\infty$, $-\dfrac{3}{2}$ ,  0 , $\dfrac{3}{2}$,  $+\infty$}
       \tkzTabLine
     		{ ,+,z,-,d,-,z,+, }
       \tkzTabVar
     		{ -/ , +/ $-11$, -D+ / / ,-/ $13$, +/   }
\end{tikzpicture}
\end{center}
\end{enumerate}
\end{corrige}

\begin{corrige}
\begin{enumerate}

\item $g'(x)=1+\dfrac{1}{x^2}$
\item D'une part on sait que $g'(x)=1+\dfrac{1}{x^2}=\dfrac{x^2+1}{x^2}$.
\item On sait que $x^2$ et $x^2+1$ sont des quantités strictement positive sur $\mathbb{R}^*$ donc on en déduit le tableau récapitulatif d'étude de la fonction $g$,
\begin{center}
\begin{tikzpicture}
     \tkzTabInit[lgt=3,espcl=3]
       {$x$ / 1 , Signe de $g'(x)$ / 1, Variation de  $g$ / 2 }
       {$-\infty$,  0 ,  $+\infty$}
       \tkzTabLine
     		{ ,+,d,+, }
       \tkzTabVar
     		{ -/ , +D- / /, +/   }
\end{tikzpicture}
\end{center}
\end{enumerate}
\end{corrige}

\begin{corrige}
\begin{enumerate}

\item $h'(x)=3x-7+\dfrac{4}{x^2}$
\item D'une part on sait que $h'(x)=3x-7+\dfrac{4}{x^2}=\dfrac{3x^3 -7x^2 +4}{x^2}$. D'autre part, 

$$(x-2)(x-1)(3x+2)=(x^2-3x+2)(3x+2)=3x^3+2x^2-9x^2-6x+6x+4=3x^3 -7x^2 +4$$
 
\item De la forme factorisée, on en déduit le tableau récapitulatif d'étude de la fonction $h$,
\begin{center}
\begin{tikzpicture}
     \tkzTabInit[lgt=3,espcl=2]
       {$x$ / 1 , Signe de $h'(x)$ / 1, Variation de  $h$ / 2 }
       {$-\infty$, -$\dfrac{2}{3}$ , 0 , 1, 2,  $+\infty$}
       \tkzTabLine
     		{ , -, z, +, d, +, z,-,z,+, }
       \tkzTabVar
     		{ +/ , -/ $\frac{34}{3}$ , +D- / /, +/ ${-9,5}$ , -/ $-10$, +/  }
\end{tikzpicture}
\end{center}
\end{enumerate}
\end{corrige}

\newpage

\begin{corrige}
\begin{enumerate}
\item $f'(x)=1- \dfrac{6400}{x^2}$
\item $$f'(x)=1- \dfrac{6400}{x^2}= \dfrac{x^2-6400}{x^2}=\dfrac{(x-80)(x+80)}{x^2}$$ 
Par suite on déduit que $f'(x)$ est du même signe que $(x-80)(x+80)$.
\item Tableau récapitulatif de la fonction $f$ 
\begin{center}
\begin{tikzpicture}
     \tkzTabInit[lgt=3,espcl=3]
       {$x$ / 1 , Signe de $f'(x)$ / 1, Variation de  $f$ / 2 }
       {$40$ , $80$,  $160$}
       \tkzTabLine
     		{ ,-, z, +, }
       \tkzTabVar
     		{ +/ $100$, -/ $60$, +/ $100$}
\end{tikzpicture}
\end{center}
\item 
\begin{enumerate}
\item Tableau à compléter 
\begin{center}
\begin{tabular}{|c|c|c|c|}
\hline 
Nombre de repas & 40 & 50 & 100 \\ 
\hline 
Coût de $x$ repas & 4 000 & 3 900 & 6 400 \\ 
\hline 
 Coût moyen d'un repas & 100 & 78 & 64 \\ 
\hline 
\end{tabular} 
\end{center}
\item $C_m(x)=\dfrac{C(x)}{x}$
\item Pour que le coût d'un repas soit minimal il faut servir 80 repas. 
\item A l'aide du graphique on trouve que le nombre de repas $x$ doit appartenir à l'intervalle $[44;146]$.
\end{enumerate}
\end{enumerate}
\end{corrige}

\end{document}
